\documentclass[11pt,a4paper]{article}

\usepackage[utf8]{inputenc}

\usepackage[T1]{fontenc}

\usepackage[polish]{babel}

\usepackage{lmodern}

\usepackage{geometry}

\geometry{a4paper,margin=27mm}

\usepackage{microtype}

\usepackage[hidelinks]{hyperref}

% Portfolio: MSC-AR-P02; source: prospective ARC-XR research scope.

\title{Frustracja i dekoncentracja w procedurach manualnych AR}

\author{Mateusz Pomianek}

\date{Wersja 0.3.0, 21 listopada 2026}

\begin{document}

\maketitle

\begin{abstract}

Przedmiotem opracowania jest Frustracja i dekoncentracja w procedurach manualnych AR. Czy indywidualna kalibracja poprawia rozpoznawanie deklarowanej frustracji względem estymatora ogólnego? Zestawione zostaną dwa modele rozpoznawania frustracji: model ogólny i indywidualnie kalibrowany, oceniane na oddzielnych sesjach z niezależnym samoopisem. Głównym punktem końcowym jest macro-F1 klasyfikacji frustracji na rozdzielonych czasowo sesjach. Opis ma charakter wstępny i nie przedstawia niezweryfikowanych wyników.

\end{abstract}

\section{Wprowadzenie}

Rozwój adaptacyjnych systemów XR stwarza możliwość badania interakcji człowieka z informacją przestrzenną, ale samo zastosowanie AR, VR lub modeli sztucznej inteligencji nie przesądza o przewadze poznawczej czy ergonomicznej. W rozpatrywanym zagadnieniu, Frustracja i dekoncentracja w procedurach manualnych AR, kluczowe jest rozdzielenie właściwości systemu od zmienności zachowania użytkownika. Porównanie należy zaprojektować tak, aby mierzyć jednoznaczne zjawisko, a nie tylko ogólną ocenę technologii.

\section{Cel pracy i pytanie badawcze}

Czy indywidualna kalibracja poprawia rozpoznawanie deklarowanej frustracji względem estymatora ogólnego? Celem minimalnego eksperymentu jest weryfikacja predykcji względem niezależnego samoopisu, bez przypisywania samemu klasyfikatorowi skuteczności interwencji.

\section{Zakres i odrębny wkład naukowy}
Predykcja deklarowanej frustracji w zadaniu manualnym AR, odrębna od skuteczności jakiejkolwiek interwencji emocjonalnej. Wariant minimalny ocenia personalizację między sesjami. Test kalibracji osobniczej względem nowej sesji, nie ogólna demonstracja VR z fizjologią.

\section{Hipotezy badawcze}
\paragraph{H1.} Model z kalibracją na sesji wcześniejszej osiąga wyższy macro-F1 względem deklaracji niż model globalny na niewidzianej sesji tej samej osoby.
\paragraph{H2.} Poprawa personalizacji pozostaje dodatnia przy jednakowym budżecie etykiet i pomiarze odsetka odmów klasyfikacji.
\paragraph{H3.} Względny zysk jest mniejszy przy większej zmianie warunków sesji, co ujawnia ograniczenie transferu indywidualnej kalibracji.

\section{Projekt badania i plan eksperymentów}
\subsection{A. Operacjonalizacja i kwalifikacja danych}
Predykcja deklarowanej frustracji w zadaniu manualnym AR, odrębna od skuteczności jakiejkolwiek interwencji emocjonalnej. Wariant minimalny ocenia personalizację między sesjami. Jednostka wnioskowania: uczestnik i niezależna sesja testowa; próbki czasowe z tej samej sesji nie są niezależne. Punkt końcowy i sposób oceny będą zamrożone przed próbą potwierdzającą. Należy udokumentować pochodzenie obserwacji, czas akwizycji, identyfikatory sesji, harmonogram oraz legalność użycia materiałów.

\subsection{B. Eksperyment porównawczy}
Ten sam zestaw cech i architektura modelu; porównać trening globalny z predefiniowaną liczbą etykiet personalizacyjnych bez użycia przyszłej sesji. Estymand pierwszorzędowy: różnica macro-F1 klasyfikacji frustracji między modelem indywidualnie kalibrowanym a globalnym na nowych sesjach. Przypisanie warunków, mianowniki i kryteria wyłączeń nie mogą być dobierane po przejrzeniu wyników.

\subsection{C. Kontrole i analiza odporności}
Zdefiniować niezależny samoopis, próg etykiety i odsetek klas; brak diagnozy klinicznej lub obiektywnego pomiaru emocji. Zmienność osoby i sesji, odrębny czas i światło, jakość sygnału oraz brakujące modalności. Każdą odmowę systemu, niekompletną próbę i brak sygnału należy raportować osobno, a nie po cichu usuwać z próby.

\subsection{D. Walidacja na niezależnym materiale}
Ocenę należy wykonać na uprzednio odłożonych scenach, sesjach lub uczestnikach. Uczestnikowo-klastrowy bootstrap różnicy macro-F1, analiza czułości na częstość klasy; dobór progów wyłącznie na treningu. Model i próg zatrzymać przed analizą końcową; rejestrować również brak przewagi i działania niepożądane.

\section{Estymand, analiza statystyczna i decyzje}
Główny estymand: różnica macro-F1 klasyfikacji frustracji między modelem indywidualnie kalibrowanym a globalnym na nowych sesjach. Jednostką powtarzalności jest uczestnik i niezależna sesja testowa; próbki czasowe z tej samej sesji nie są niezależne. Uczestnikowo-klastrowy bootstrap różnicy macro-F1, analiza czułości na częstość klasy; dobór progów wyłącznie na treningu. H1 jest kandydatem na hipotezę potwierdzającą. H2 i H3 pozostają eksploracyjne, dopóki protokół nie określi ich hierarchii i kontroli wielokrotnych porównań. Minimalny efekt istotny praktycznie, liczebność prób, margines nie gorszości i progi jakości wyznaczyć na podstawie pilotażu i analizy precyzji, nie zgadywać wartości.

\section{Stan badań i granica nowości}
\begin{itemize}
\item Physiology-driven VR adaptation has been demonstrated; independent behavioural gains require direct controlled trials. \cite{xrprior1}.
\item EEG/VR time alignment and experimental validity must be explicitly documented before any multimodal inference claim. \cite{xrprior2}.
\end{itemize}
Test kalibracji osobniczej względem nowej sesji, nie ogólna demonstracja VR z fizjologią. Samo zastosowanie VLM, EEG, okulometrii lub adaptacyjnego HUD nie stanowi wystarczającej nowości międzynarodowej.

\section{Warunki bezpieczeństwa, sprzętu i zgodności}
Opis stanowi prospektywny plan, a nie prerejestrację ani zatwierdzoną procedurę. Przed zbieraniem danych od ludzi niezbędna jest właściwa weryfikacja etyczna i zgoda. Standardowy Quest 3 nie ma wbudowanego eye trackingu, a Tobii Eye Tracker 5 działa na stanowisku ekranowym. Odstępstwa, obciążenia i ograniczenia czujników muszą być ujawnione. Nie publikować surowych danych biometrycznych ani nagrań uczestników.

\section{Materiał i metody}

Zestawione zostaną dwa modele rozpoznawania frustracji: model ogólny i indywidualnie kalibrowany, oceniane na oddzielnych sesjach z niezależnym samoopisem. Wymagane jest zdefiniowanie jednostki analizy, zasad doboru przypadków, wersji scenariusza i źródeł pomiaru. Rozkład warunków oraz reguły wykluczenia prób należy określić przed pomiarem. W badaniach z udziałem ludzi niezbędna jest uprzednia właściwa ocena etyczna i świadoma zgoda. W testach wyłącznie symulacyjnych należy oddzielić prawdę symulacyjną od obserwacji empirycznej.

\subsection{Wynik główny i kryteria jakości}

Macro-F1 klasyfikacji frustracji na rozdzielonych czasowo sesjach. Pomiar powinien zachowywać jawne jednostki, znaczniki czasu i reguły oceny poprawności. Przy modelach predykcyjnych stosować rozłączne partycje oceny; przy eksperymentach interwencyjnych oddzielić skuteczność predykcji od przyczynowego wpływu interfejsu.

\section{Plan analizy}

Zestawić Macro-F1 z przedziałami ufności po uczestnikach, porównywalne budżety danych oraz ocenę kalibracji predykcji. W wynikach należy raportować wielkość efektu lub błąd estymacji z przedziałem niepewności, liczbą obserwacji i odsetkiem odrzuconych prób. Analizy dodatkowe wymagają jasnego oznaczenia jako eksploracyjne. Nie wybierać progów decyzyjnych ani modelu na podstawie obserwacji ze zbioru testowego.

\section{Wyniki}

W obecnej wersji nie raportuje się wyników pomiarów, tabel statystycznych ani parametrów wyznaczonych z nieudokumentowanych danych. Nie stwierdza się w tym opracowaniu, że eksperyment został wykonany lub zgłoszony do czasopisma.

\section{Dyskusja}

Interpretacja przyszłych wyników musi uwzględnić alternatywne mechanizmy i zależność wniosków od jakości danych, warunków zadania oraz ograniczeń sprzętu. Wartość potencjalnego rozwiązania zależy nie tylko od trafności modelu, lecz także od jego opóźnienia, stabilności i wpływu na zachowanie człowieka. Wyniki testów syntetycznych nie mogą być automatycznie uogólniane na populację uczestników.

\section{Ograniczenia badania}

Samoopis nie jest bezpośrednią prawdą neuronalną; należy zapobiec przeciekowi danych między treningiem a testem. Granice nowości i trafności ekologicznej wymagają niezależnej weryfikacji. Żadna część szkicu nie dowodzi przeprowadzenia zbierania danych, uzyskania zgód etycznych ani gotowości manuskryptu do zgłoszenia.

\section{Wnioski}

Przedstawiono problem, kontrast i sposób jego planowanej oceny. Rozstrzygnięcie postawionego pytania wymaga przygotowania protokołu, zweryfikowanych źródeł oraz poprawnie udokumentowanych wyników. Bez tych materiałów nie można sformułować konkluzji empirycznych.

\begin{thebibliography}{9}
\bibitem{xrprior1} Towards Intelligent VR Training: A Physiological Adaptation Framework for Cognitive Load and Stress Detection (2025), UMAP 2025. \url{https://doi.org/10.1145/3699682.3727575}.
\bibitem{xrprior2} The integration of virtual reality and EEG: A step-by-step guideline (2026), Methods guideline. \url{https://pmc.ncbi.nlm.nih.gov/articles/PMC12808493/}.
\end{thebibliography}
\end{document}

